\chapter{相关技术}

\section{视觉语言预训练模型 CLIP}

CLIP 是当前视觉语言预训练研究中最具代表性的模型之一\cite{clip2021}。它的核心思想是通过大规模图文对比学习,把图像和文本映射到同一语义空间中。训练时,模型希望匹配样本的图像特征与文本特征更接近,不匹配样本则更远。经过这种方式训练后,视觉编码器学到的是与语义概念相关的视觉表示,文本编码器学到的是能够与视觉语义对齐的语言表示。

从结构上看,CLIP 由视觉编码器和文本编码器两个分支组成。视觉编码器可以采用 ResNet 或 Vision Transformer,文本编码器则通常基于 Transformer 架构\cite{clip2021,chen2023vlp}。图像经过视觉编码器得到视觉特征,文本模板经过分词和编码后得到文本特征,二者再通过余弦相似度完成匹配。对于分类任务,只需把类别名称填入文本模板,即可把文本特征视为分类权重。

CLIP 的优势在于它不依赖任务专用分类头,理论上可以借助自然语言直接描述新类别,从而具备较强的零样本迁移能力\cite{clip2021}。对于标注数据较少的细粒度任务,这种能力尤其有吸引力。不过,CLIP 在下游任务中的表现对提示模板高度敏感。对于类别差异细微的任务而言,如果文本模板过于通用,文本语义往往不足以支撑不同子类之间的有效区分。

\section{提示学习方法}

\subsection{CoOp}

CoOp 是将提示学习系统性引入视觉语言模型适配的代表方法\cite{coop2022}。传统 CLIP 往往使用人工撰写的固定模板,如 ``a photo of a \{class\}''。CoOp 的做法是保留类别名称,把模板中的上下文词替换为可学习向量,再通过少量标注样本优化这些向量。

设提示长度为 \(n_{ctx}\),每个上下文 token 的维度与文本编码器嵌入维度一致,则 CoOp 实际学习的是一组连续上下文表示
\begin{equation}
\mathbf{C}=\{\mathbf{c}_1,\mathbf{c}_2,\ldots,\mathbf{c}_{n_{ctx}}\}.
\end{equation}
这些向量与类别名称 token 拼接后送入文本编码器,得到新的文本特征。训练时通常冻结视觉编码器和文本编码器主干,只更新提示参数,因此额外训练成本较低。

CoOp 的意义在于证明了少量可学习提示参数就能显著提升 CLIP 在下游分类任务上的表现\cite{coop2022}。这说明视觉语言模型的任务适配并不一定要依赖全模型微调,提示本身就能承载相当一部分任务信息。

但 CoOp 的提示是静态共享的,所有输入图像使用同一组上下文向量。当样本分布复杂或类别边界模糊时,这种统一提示很难同时适应所有样本。对于细粒度分类,这一局限更明显,因为不同图像需要关注的局部差异并不相同。

\subsection{CoCoOp}

为了解决静态提示的限制,Zhou 等人进一步提出 CoCoOp\cite{cocoop2022}。该方法在 CoOp 基础上增加图像条件机制,通过一个轻量级网络根据当前图像特征生成条件 token,再与原有提示结合,形成面向当前输入的动态提示。

如果把图像特征记为 \(\mathbf{f}_v\),则 CoCoOp 可抽象为
\begin{equation}
\mathbf{t}_{cond}=g(\mathbf{f}_v),
\end{equation}
其中 \(g(\cdot)\) 是一个小型神经网络。生成的条件 token 会影响文本提示,从而让不同图像拥有不同的上下文偏移。

与 CoOp 相比,CoCoOp 的价值在于它把提示学习从“任务级共享”推进到“样本级适配”。同一类别下的不同图像,可能因为姿态、遮挡、背景和局部纹理差异而需要不同的提示表达; 条件提示正是为了解决这个问题而设计的。原论文表明,这一机制有助于提升模型在少样本泛化和跨数据集测试中的表现\cite{cocoop2022}。

不过,CoCoOp 主要关注“如何根据图像生成提示”,并没有进一步回答“训练时哪些样本更应该被重点学习”。也就是说,它虽然提高了提示的动态性,但没有改变损失函数对不同训练样本一视同仁的处理方式。

\subsection{自动提示优化与多模态提示}

在 CoOp 和 CoCoOp 之后,提示学习开始向更复杂的方向扩展。一类工作关注自动提示优化。ProAPO 通过渐进式策略对提示进行自动搜索和改进,尽量减少人工模板设计对结果的影响\cite{qu2025proapo}。这类方法说明,提示优化并不局限于简单的梯度更新,还可以结合搜索和逐步筛选机制。

另一类工作强调多模态提示协同。传统提示学习多集中在文本侧,而多模态提示方法会同时考虑视觉提示和文本提示之间的交互。Liu 等人在 Neurocomputing 2025 中提出细粒度多模态提示学习方法,指出在视觉语言模型中同时调节视觉提示与文本提示,有助于增强对细粒度差异的表达能力\cite{liu2025fgmpl}。类似地,更一般的多模态提示优化研究也表明,如果只调整单一模态提示,往往无法充分利用模型内部已有的跨模态对齐能力\cite{choi2025multimodal}。

\section{细粒度图像分类相关方法}

细粒度分类的核心难点在于如何稳定提取具有区分性的局部特征。Oxford-IIIT Pets 的提出就明确说明,宠物品种识别是一类具有高变形、高姿态变化和细微类别差异的问题\cite{parkhi2012cats}。在这种任务中,全局语义通常只能帮助模型判断大类,真正决定分类结果的往往是局部毛色、耳型、脸部轮廓和体态比例。

在纯视觉路线中,卷积网络和迁移学习仍然是常见方案。针对犬种识别,已有研究采用卷积网络结合传统分类器进行实验,结果表明类别相似性和数据复杂度仍然是限制性能的重要因素\cite{cuizhang2023classification}。在猫类识别场景下,基于迁移学习的分层特征利用也被证明能够改善细粒度分类表现\cite{wang2024enhancing}。这些研究说明,即使不考虑视觉语言模型,细粒度宠物分类本身就具有较高难度。

围绕小样本适配,视觉 Transformer 相关研究也开始强调轻量化参数更新的重要性。例如,针对少样本适配的研究指出,在冻结大部分 backbone 参数的前提下进行少量参数调节,通常比直接全量微调更稳定\cite{xu2023fewshotvit}。这种思路与提示学习的出发点本质一致,即尽量通过少量可训练参数完成任务适配。

从更大的研究背景来看,细粒度分类与视觉语言提示学习的结合正在成为一个逐渐清晰的方向。一方面,零样本和少样本设置能够放大视觉语言模型的预训练优势; 另一方面,细粒度类别差异足够微弱,也能更直接地检验提示设计是否真正提高了类别区分能力\cite{atabuzzaman2025zero,liu2025fgmpl}。因此,在 CLIP 基础上研究细粒度猫狗分类,既有方法研究意义,也有稳定的数据基础。

\section{难样本学习与损失加权}

除了动态提示外,本文方法还强调对困难样本的关注。因此,有必要回顾难样本学习相关思路。

标准交叉熵默认所有样本贡献相同,这在很多任务中已经足够有效。但对于类别边界复杂或难例比例较高的任务,统一权重可能使模型更多地被简单样本主导。Focal Loss 就是解决这一问题的经典方法之一\cite{lin2017focal}。它通过降低易分类样本的损失权重,让模型把更多注意力放到难分类样本上。其形式为
\begin{equation}
\mathrm{FL}(p_t) = -\alpha_t (1-p_t)^\gamma \log(p_t),
\end{equation}
其中 \(p_t\) 表示真实类别预测概率,\(\gamma\) 用于调节难样本聚焦强度。

虽然 Focal Loss 最早主要用于目标检测,但其核心思想同样适用于细粒度分类。对于本文任务而言,真正决定模型性能的往往是那些外观非常接近、预测置信度较低的品种样本。因此,在提示学习框架中引入基于预测结果的样本加权,是一个合理且具有针对性的改进方向。

\section{本章小结}

本章从四个方面介绍了与本文工作相关的技术基础。首先,回顾了 CLIP 的核心结构与优势,说明视觉语言预训练为少样本细粒度分类提供了有力的表示基础。其次,梳理了 CoOp、CoCoOp 以及自动提示优化、多模态提示等提示学习方法,明确了静态提示、条件提示和更复杂提示策略之间的差异。然后,结合宠物品种识别相关研究说明细粒度分类本身的任务特点与挑战。最后,讨论了难样本学习和损失加权的基本思路,为本文在提示学习中引入难度感知权重机制提供了理论依据。
